\section{Introduction}\label{sec:intro}

Tests of the mean of a stationary series with long memory require a critical value that depends on the Hurst index
$H$ \citep{Ber94,BFGK13}. We study three ways of normalising the mean: the self-normaliser of \citet{Lob01,Sha10,Sha15},
equivalent to the Bartlett kernel at full bandwidth \citep{KVB00,KV02,KV05,MP12}; the fixed-$m$ periodogram
statistic of \citet{HI17}; and the equal-weighted cosine (EWC) statistic \citep{Mul14,Sun14,LLSW18}.
For each, we ask how accurately its limit critical values can be computed, how often the resulting test rejects
in finite Gaussian samples and how its behaviour changes when $H$ must be estimated.

At $H=\frac12$, the self-normalised law is known in closed form \citep{AP97,AP02,Lob01,KV02}. At other values of $H$,
the tables of \citet{MP09,MP12,MP13,HI17} rely on simulation. Comparing their entries requires matching the tail,
the signed or squared statistic and its normalisation. Expressing numerical discrepancies in rejection-probability
units gives these comparisons a common scale.

For the Bartlett limit $T_H=B_H(1)^2/\int_0^1\Bt(r)^2\dd r$, we enclose the truncation error of a cosine expansion
of the fractional Brownian bridge and integrate the resulting Gaussian quadratic-form law in ball arithmetic.
The fixed-smoothing limits admit a finite-dimensional calculation for $0<H<1$.
For Gaussian samples with specified covariance, the same quadratic-form representation gives finite-sample
rejection probabilities. We use Monte Carlo to check these probabilities and to evaluate feasible procedures
that estimate $H$ or use subsampling.

\paragraph{Contributions}
\begin{enumerate}
\item[(C1)] Certified critical values of $T_H$ at six levels and of the fixed-$m$ and EWC limits at three levels, tabulated
for $\frac12\le H\le0.99$, with deterministic enclosures (\cref{sec:cv,sec:computation}, \cref{thm:certificate}) and adaptive
quadrature in complex ball arithmetic \citep{Joh18}. An analytical bound on the far tail speeds up assembly.
On benchmark T4 ($H=0.75$, $J=16384$), one certified probability takes \pthree{e5-T4-cert-s}\,s including setup.
The fixed-smoothing certificate also applies for $0<H<\frac12$. As $H\uparrow1$, fixed quantiles of the squared
statistics grow at rate $(1-H)^{-1}$ and their positive signed cutoffs at rate $(1-H)^{-1/2}$ (\cref{prop:edge}).
\item[(C2)] Exact finite-sample rejection rates under Gaussian long-memory models via a Cholesky-free Woodbury
eigendecomposition (\cref{prop:exactn}), validated against independent Monte Carlo, together with a head-to-head
feasible size and local-power comparison against exact local Whittle plug-in, conservative upper-confidence plug-in,
MAC variance estimation \citep{ADG09} and rate-free subsampling \citep{PRW99} (\cref{sec:finite}).
\item[(C3)] A closed-form first-moment rescaling of the Student-$t$ values for the fixed-smoothing tests with exactly
computed size error, and the collapse of the self-normaliser that explains why it works and why it fails for the
Bartlett/self-normalised statistic (\cref{prop:rho}).
\item[(C4)] Local power, exact for the fixed-smoothing tests and computed for the Bartlett/self-normalised test, and the
size--power trade-off across the three families (\cref{sec:power}).
\item[(C5)] Empirical demonstrations on three contiguous series from energy systems, urban air quality and hydrology
(\cref{sec:empirical}), quantifying decision reversals and confidence-interval expansion in practice.
\item[(C6)] An audit, in test-size units, of \pthree{audit-n-values} published critical values from
\pthree{audit-n-sources} sources, with explicit matching of statistic, normalisation, tail and design
(\cref{sec:audit}, Supplementary Material Section~S.3).
\item[(C7)] A benchmark of \pthree{e5-n-impl} implementations of \pthree{e5-n-methods} computational methods, including
four modes of the software of \citet{Das25}, on identical targets, charging each method its own spectral step and
reporting error bounds against certified references where available (\cref{sec:methods}).
\end{enumerate}
The mathematical tools build on the companion paper \citep{AhmedCosChaos}: closed-form cosine Gram matrices of
fractional Brownian motion and a two-sided weighted cosine inequality. Here they yield a two-sided bound on every
eigenvalue of the fractional Brownian bridge covariance operator (\cref{thm:eig}), which provides the remainder estimate
the certificate needs. We call the resulting certified computation \CosCert{}, for cosine-coordinate certificate; its
reproducible code accompanies the paper. The classical case $H=\frac12$ \citep{AP97} validates the engine to
\pthree{g0-min-digits}--\pthree{g0-max-digits} significant digits.

\paragraph{Organisation} \Cref{sec:setting} fixes the statistics, their normalisations and their common coordinates.
\Cref{sec:cv} gives the critical values and response surface, \cref{sec:finite} the finite-sample rejection rates and
feasible competitor comparison, \cref{sec:power} local power and test selection, \cref{sec:empirical} the three
empirical applications, \cref{sec:audit} the audit of published tables, \cref{sec:computation} the \CosCert{}
certificate and computational benchmark, and \cref{sec:conclusion} practical recommendations. Proofs are in
\cref{app:proofs} and certificate details in \cref{app:certificate}; complete critical-value tables, extended simulation
tables, the full 233-value audit ledger and validation details are provided in the standalone Supplementary Material
(Sections~S.1--S.4).
